% 按原书第 121 页重排表 3.5。
\begin{table}[p]
\centering
\makebox[\linewidth][c]{\rotatebox{90}{%
\begin{minipage}{.94\textheight}
\caption{多提示推理方法。}
\label{tab:ch03-multi}
\footnotesize
\setlength{\tabcolsep}{3pt}
\renewcommand{\arraystretch}{1.14}
\begin{tabularx}{\linewidth}{@{}>{\raggedright\arraybackslash}p{.095\linewidth}
  >{\raggedright\arraybackslash}p{.173\linewidth}
  >{\raggedright\arraybackslash}X
  >{\raggedright\arraybackslash}p{.16\linewidth}
  >{\raggedright\arraybackslash}p{.17\linewidth}@{}}
\toprule
方法 & 概述 & 示例 & 优点 & 缺点\\
\midrule
集成 &
用多个模板完成一项推理任务，以简单多数投票、平均、蒸馏等办法，将每个测试样例的各次预测汇总成最终答案。 &
提示 1\enspace “I am in the capital of Canada, I am in [MASK] .”\par
提示 2\enspace “[MASK] is located in Canada, and is its capital city .”\par
提示 3\enspace “[MASK] is the capital of Canada”\par
各提示的首选结果\enspace 1) Toronto，2) Ottawa，3) Ottawa\par
简单多数投票\enspace Ottawa（正确） &
对多个提示求平均，可以减轻个别提示失效带来的影响。\par
可充分利用提示工程方法生成的多个候选模板。 &
推理时，每个测试样例需要多次预测。\par
效果受输入提示的质量限制，优化余地较小。\\
\midrule
上下文学习 &
推理前，在尚未填入答案的模板之前或之后添加已填好查询—答案对的模板示例，引导大语言模型给出所需的预测。 &
原提示\enspace “The capital of Canada is [MASK] .”\par
上下文提示\enspace “The capital of France is Paris. The capital of Canada is [MASK] .”\par
两种提示的首选结果\enspace 1) Ottawa（$p=0.284$），2) Ottawa（$p=0.341$） &
推理时，每个测试样例只需预测一次。\par
实现简单、直观。 &
准确率很大程度上取决于所选的上下文示例。\par
要得到好的效果，可能需要对上下文提示微调。\\
\midrule
提示分解 &
构造多个子提示，分别探查复杂查询的不同方面，例如识别一句话中的各个命名实体。 &
原提示\enspace \texttt{[x]} = “Serena won the tennis tournament”\par
分解后的提示\par
“[x] . ‘Serena’ is a [MASK]”（人名）\par
“[x] . ‘Tennis’ is a [MASK]”（运动）\par
“[x] . ‘the’ is a [MASK]”（非实体） &
可探查查询中粒度更细的细节。 &
推理时，每个测试样例需要多次预测。\par
可能需要逐一优化多个子提示。\\
\bottomrule
\end{tabularx}
\end{minipage}}}
\end{table}
